\section{Integration provenance and evidence boundaries}\label{sec:integration}

\subsection{What was consolidated}

The incoming expert responses N1--N9 answered the same 31 questions in the
earlier version of this article. They are not the older architecture proposals
P1--P9 retained in directories 01--09. Their common explanations now have one
maintained home: the semantic boundary in
Section~\ref{sec:semantic-boundaries}, the topic-specific answers in Part~III,
and the roadmap in Section~\ref{sec:roadmap}. Distinct constructions and
assumptions are retained; repeating the same claim in several responses is
not additional experimental evidence.

The complete original packages are pinned to repository revision
\code{5c3a53c} (the full SHA is recorded in the inventory). The machine-readable
\path{integration/sources.json} identifies all 97 original files by path,
Git blob and SHA-256 and records their disposition. Original articles and PDFs
remain recoverable from that revision; they are superseded as reading paths
by this maintained article. The 28 distinct executable-evidence and provenance
files are preserved byte-for-byte under \path{evidence/N1/} through
\path{evidence/N9/}. Source READMEs and standalone article build instructions
are superseded by the maintained reproduction instructions, not treated as
additional design recommendations.

\begin{longtable}{@{}L{10mm}L{71mm}L{68mm}@{}}
\toprule ID & Original package/article & Distinct evidence emphasis\\\midrule\endhead
N1 & \path{new/1/article.tex}, split sections & Binary length-modulo-three carrier, conflict clique, eleven model checks.\\
N2 & \path{new/2/Leant2_Expert_Answers/} & Four-state two-symbol machine; direct-tail obstruction and finite search accounting.\\
N3 & \path{new/3/leant2-expert-answers/} & Unary modulo-three, rollback reuse and approximation-direction countermodels.\\
N4 & \path{new/4/leant2_expert_answers/} & Parity and unary carriers, conjunction-aware retrieval, partial-row graph.\\
N5 & \path{new/5/leant2_expert_answers/} & Four-state witness, strict request/publication distinctions and negative-policy audit recommendations.\\
N6 & \path{new/6/leant2_expert_answers.tex} & Option min/max, Fibonacci pair, observational noncongruence and branch-qualified cache model.\\
N7 & \path{new/7/Leant2_Expert_Questions/} & Feature-cover enumeration, shared angelic-call conflict, historical 31-question index.\\
N8 & \path{new/8/leant2_expert_answers/} & Minimal sample-consistent carrier that overfits and is refined by a new observation.\\
N9 & \path{new/9/leant2-expert-answers.tex} & Maximum-prefix-sum invariant monoid, scheme-dependent state and function-carrier reversal.\\
\bottomrule
\end{longtable}

These path names are provenance identifiers, not active links to the retired
incoming directories. The integration ledgers next to the source inventory
map source sections and individual claims to their canonical destinations.
They record duplicate treatment, corrections, alternatives and question
dispositions. The bibliography consolidates references by work and pinned
source rather than preserving multiple aliases for the same publication.

\subsection{The historical question inventory}

Answered questions are removed from the active expert agenda. The following
table preserves their identifiers and scope for provenance without posing them
again. Each topic's disposition paragraph records the established answer and
the remaining limitation. New questions use the distinct suffix \code{N};
they ask about the residual problem, not whether an already identified method
exists.

\begin{longtable}{@{}L{22mm}L{87mm}L{38mm}@{}}
\toprule Historical IDs & Themes now covered by integrated answers & Canonical section\\\midrule\endhead
R1.Q1--Q4 & Coupled goals, admissible cost bounds, learning versus ranking, deadline reproducibility & \ref{sec:search}\\
R2.Q1--Q5 & Carrier minimality, polymorphic alphabets, fusion/lifting, motive bases, generalization & \ref{sec:carriers}\\
R3.Q1--Q4 & Higher-order live evaluation, rollback-safe normalization, decisive evaluation, bounded angelic search & \ref{sec:evaluation}\\
R4.Q1--Q4 & Motive templates, transports, unfolding, native induction interfaces & \ref{sec:dependent}\\
R5.Q1--Q3 & Interpreter agreement, top-down recursion constraints, practical recursion schemes & \ref{sec:recursion}\\
R6.Q1--Q4 & Value/proof interaction, sufficient induction fragments, automation, proof scheduling & \ref{sec:contracts}\\
R7.Q1--Q3 & Dependent abstraction, type-class prerequisites, relevance for programs & \ref{sec:retrieval}\\
R8.Q1--Q2 & Ranking objectives and exact versus sampled equivalence & \ref{sec:ranking}\\
R9.Q1--Q2 & Model-proposal evaluation and replay/auditability & \ref{sec:learning}\\
\bottomrule
\end{longtable}

Several answers are conditional rather than universal: finite-state minimality
does not identify a minimal intended Lean program; a finite motive grammar
does not enumerate all dependent motives; a measurement protocol does not
supply an unmeasured model-accuracy percentage. Retiring an overbroad or
answered question therefore does not assert that all associated research is
finished. The active agenda in Section~\ref{sec:experts} points to the new,
more specific questions.

\subsection{What the preserved executable checks establish}

All nine Python suites were rerun in isolated copies on September 22, 2026.
Every process exited zero with empty stderr; every decoded JSON result matched
its archived result; all preserved input hashes were unchanged. The raw
process streams, executable hash and receipt are under
\path{evidence/rerun-2026-09-22/}. The driver
\path{tools/check_proposal_models.py} recreates the experiment in a fresh
directory and never overwrites the historical receipts.

The suites have different mathematical domains and stopping rules. N1's
binary machine search exhausts the one- and two-state spaces and stops on a
three-state witness. N2/N7 account for different seed conventions than N5;
their table counts must not be interchanged. N3/N6/N9 use unary modulo-three
models. N8 stops on a minimal four-state sample witness which fails at length
six, then refines it. N9 checks maximum-prefix-sum closure and associativity
on a finite invariant domain, including a counterexample to unrestricted
identity. Written general proofs and the finite executions remain distinct.

These checks execute neither Lean nor Leant2. They establish no native closure
semantics, motive elaboration, engine throughput, model accuracy or universal
refutation theorem. The original source manifests and checksum reports also
retain their historical meaning: their article page counts and file hashes
refer to the source packages, not to this consolidated PDF. The documentation
build and visual review are recorded separately.

\subsection{Versioned source observations and disagreements}

The original performance table belongs to \code{33cec8b}; the responses'
selected source inspections belong to \code{784664f}. Pinned Lean 4.34.0
source advice is distinguished from moving manual pages consulted by the
authors, some describing 4.35.0-rc2. A historical source observation is not a
fresh implementation test. Reported external corpus counts and cloud timing
conditions remain measurements of those external systems.

The consolidation does not turn differing confidence levels into agreement:
documented \code{sauto} options do not prove that its search policy transfers
to Lean; native motive construction does not solve unrestricted generalization;
equation-based evaluation need not coincide with definitional reduction;
library selection is not a complete dependent retrieval procedure. The
topic sections retain these qualifications where they affect an algorithm.
The six independent obligations of
Section~\ref{sec:semantic-boundaries} govern all resulting claims.
